\numoptsfalse
\yearpaper{2017}{60 Questions \textperiodcentered{} 60 Marks \textperiodcentered{} 80 Minutes}{}
\begin{multicols}{2}

\QQ{A substance of mass $49.53\unit{g}$ occupies $1.5\unit{cm^3}$ of volume. The density of the substance (in $\mathrm{g\,cm^{-3}}$) with correct number of significant figures is}
{\opt{3.302}{3.300}{3.3}{3.30}}

\QQ{A car moving with a velocity of $20\unit{m\,s^{-1}}$ is stopped in a distance of $40\unit{m}$. If the same car is travelling at double the velocity, the distance travelled by it for same retardation is}
{\opt{$640\unit{m}$}{$320\unit{m}$}{$1280\unit{m}$}{$160\unit{m}$}}

\QQ{The angle between velocity and acceleration of a particle describing uniform circular motion is}
{\opt{$45^\circ$}{$60^\circ$}{$90^\circ$}{$180^\circ$}}

\QQ{If $\vec A = 2\hat i+3\hat j+8\hat k$ is perpendicular to $\vec B = 4\hat j-4\hat i+\alpha\hat k$, then the value of `$\alpha$' is}
{\opt{$\dfrac12$}{$-\dfrac12$}{1}{$-1$}}

\QQ{A body of mass $50\unit{kg}$ is suspended using a spring balance inside a lift at rest. If the lift starts falling freely, the reading of the spring balance is}
{\opt{$= 50\unit{kg}$}{$> 50\unit{kg}$}{$< 50\unit{kg}$}{$= 0$}}

\QQ{A motor pump lifts 6 tonnes of water from a well of depth $25\unit{m}$ to the first floor of height $35\unit{m}$ from the ground floor in 20 minutes. The power of the pump (in kW) is [$g = 10\unit{m\,s^{-2}}$]}
{\opt{3}{6}{1.5}{12}}

\QQ{Two balls are thrown simultaneously in air. The acceleration of the centre of mass of the two balls when in air,}
{\opt{depends on the masses of the two balls}{depends on the speeds of the two balls}{is equal to $g$ (Acceleration due to gravity)}{depends on the direction of motion of the two balls.}}

\QQ{The value of acceleration due to gravity at a depth of $1600\unit{km}$ is equal to [Radius of earth $= 6400\unit{km}$]}
{\opt{$9.8\unit{m\,s^{-2}}$}{$4.9\unit{m\,s^{-2}}$}{$19.6\unit{m\,s^{-2}}$}{$7.35\unit{m\,s^{-2}}$}}

\QQ{`Young's modulus' is defined as the ratio of}
{\opt{tensile stress and longitudinal strain}{hydraulic stress and hydraulic strain}{shearing stress and shearing strain}{bulk stress and longitudinal strain}}

\QQ{`Hydraulic lift' works on the basis of}
{\opt{Stoke's law}{Toricelli's law}{Pascal's Law}{Bernoulli's Law}}

\QQ{The S.I.\ unit of specific heat capacity is}
{\opt{$\mathrm{J\,mol^{-1}K^{-1}}$}{$\mathrm{J\,kg^{-1}K^{-1}}$}{$\mathrm{J\,K^{-1}}$}{$\mathrm{J\,kg^{-1}}$}}

\QQ{For which combination of working temperatures, the efficiency of `Carnot's engine' is the least?}
{\opt{$60\unit{K}$, $40\unit{K}$}{$40\unit{K}$, $20\unit{K}$}{$80\unit{K}$, $60\unit{K}$}{$100\unit{K}$, $80\unit{K}$}}

\QQ{The mean energy of a molecule of an ideal gas is}
{\opt{$2KT$}{$\dfrac32KT$}{$KT$}{$\dfrac12KT$}}

\QQ{Two simple pendulums A and B are made to oscillate simultaneously and it is found that A completes 10 oscillations in $20\unit{sec}$ and B completes 8 oscillations in $10\unit{sec}$. The ratio of the lengths of A and B is}
{\opt{$\dfrac85$}{$\dfrac{64}{25}$}{$\dfrac54$}{$\dfrac{25}{64}$}}

\QQ{The waves set up in a closed pipe are}
{\opt{Transverse and Progressive}{Longitudinal and Stationary}{Transverse and Stationary}{Longitudinal and Progressive}}

\QQ{Two spheres of electric charges $+2\unit{nC}$ and $-8\unit{nC}$ are placed at a distance `$d$' apart. If they are allowed to touch each other, what is the new distance between them to get a repulsive force of same magnitude as before?}
{\opt{$\dfrac{4d}3$}{$\dfrac{3d}4$}{$d$}{$\dfrac d2$}}

\QQ{Three point charges of $+2q$, $+2q$ and $-4q$ are placed at the corners A, B and C of an equilateral triangle ABC of side `$x$'. The magnitude of the electric dipole moment of this system is}
{\opt{$2qx$}{$2\sqrt3\,qx$}{$3\sqrt2\,qx$}{$3qx$}}

\QQ{$4\times10^{10}$ electrons are removed from a neutral metal sphere of diameter $20\unit{cm}$ placed in air. The magnitude of the electric field (in $\mathrm{N\,C^{-1}}$) at a distance of $20\unit{cm}$ from its center is}
{\opt{5760}{1440}{640}{Zero}}

\QQ{Two point charges $A = +3\unit{nC}$ and $B = +1\unit{nC}$ are placed $5\unit{cm}$ apart in air. The work done to move charge B towards A by $1\unit{cm}$ is}
{\opt{$1.35\times10^{-7}\unit{J}$}{$2.7\times10^{-7}\unit{J}$}{$2.0\times10^{-7}\unit{J}$}{$12.1\times10^{-7}\unit{J}$}}

\QQ{A system of 2 capacitors of capacitance $2\unit{\mu F}$ and $4\unit{\mu F}$ is connected in series across a potential difference of $6\unit{V}$. The electric charge and energy stored in the system are}
{\opt{$10\unit{\mu C}$ and $30\unit{\mu J}$}{$36\unit{\mu C}$ and $108\unit{\mu J}$}{$8\unit{\mu C}$ and $24\unit{\mu J}$}{$1\unit{\mu C}$ and $3\unit{\mu J}$}}

\QQ{The minimum value of effective capacitance that can be obtained by combining 3 capacitors of capacitances $1\unit{pF}$, $2\unit{pF}$ and $4\unit{pF}$ is}
{\opt{$\dfrac47\unit{pF}$}{$1\unit{pF}$}{$\dfrac74\unit{pF}$}{$2\unit{pF}$}}

\QQ{A cylindrical conductor of diameter $0.1\unit{mm}$ carries a current of $90\unit{mA}$. The current density (in $\mathrm{A\,m^{-2}}$) is ($\pi\simeq3$)}
{\opt{$1.2\times10^7$}{$2.4\times10^7$}{$3\times10^6$}{$6\times10^6$}}

\QQ{A piece of copper is to be shaped into a conducting wire of maximum resistance. The suitable length and diameter are \rule{1cm}{0.4pt} and \rule{1cm}{0.4pt} respectively.}
{\opt{L and d}{2L and d}{L/2 and 2d}{2L and d/2}}

\QQ{Of the following graphs, the one that correctly represents the I-V characteristics of a `Ohmic device' is}
{\optfig{\gaxes{$V$}{$I$}{\draw (0.2,0) .. controls (0.5,1.2) and (1.2,1.2) .. (1.5,0);}}
{\gaxes{$V$}{$I$}{\draw (0.1,0.1) .. controls (0.9,0.2) and (1.4,0.6) .. (1.6,1.5);}}
{\gaxes{$V$}{$I$}{\draw (0,0) -- (1.5,1.4);}}
{\gaxes{$V$}{$I$}{\draw (0.3,1.6) .. controls (0.4,0.5) and (0.9,0.3) .. (1.7,0.3);}}}

\QQ{The value of I in the figure shown below is
\fig{\begin{tikzpicture}[scale=0.75]
\draw (0,0) rectangle (2,2);
\draw[->] (-0.6,2.6) -- (0,2) node[pos=0,left,lbl]{20 A};
\draw[->] (2.6,2.6) -- (2,2) node[pos=0,right,lbl]{4 A};
\draw[->] (0,0) -- (-0.6,-0.6) node[left,lbl]{3 A};
\draw[->] (2,0) -- (2.6,-0.6) node[right,lbl]{I};
\draw[->] (0,1.2) -- (0,0.8) node[right,lbl]{5 A};
\end{tikzpicture}}}
{\opt{8\,A}{21\,A}{19\,A}{4\,A}}

\QQ{The power dissipated in $3\unit{\Omega}$ resistance in the following circuit is
\fig{\begin{circuitikz}[scale=0.55,transform shape,font=\Large]
\draw (0,0) circle (2.4);
\draw (150:2.4) to[battery1,l={4.5 V, $1\,\Omega$}] (115:2.4);
\draw (190:2.4) to[R,l_=$2\,\Omega$] (0,-0.2) to[R,l_=$6\,\Omega$] (10:2.4);
\draw (0,-0.2) to[R,l=$3\,\Omega$] (80:2.4);
\draw (-30:2.4) to[R,l_=$4\,\Omega$] (-80:2.4);
\end{circuitikz}}}
{\opt{$0.75\unit{W}$}{$0.25\unit{W}$}{$1\unit{W}$}{$0.5\unit{W}$}}

\QQ{In metre bridge experiment, with a standard resistance in the right gap and a resistance coil dipped in water (in a beaker) in the left gap, the balancing length obtained is `$l$'. If the temperature of water is increased, the new balancing length is}
{\opt{$>l$}{$<l$}{$=l$}{$=0$}}

\QQ{A proton, a deuteron and an $\alpha$-particle are projected perpendicular to the direction of a uniform magnetic field with same kinetic energy. The ratio of the radii of the circular paths described by them is}
{\opt{$1:\sqrt2:1$}{$1:\sqrt2:\sqrt2$}{$\sqrt2:1:1$}{$\sqrt2:\sqrt2:1$}}

\QQ{A galvanometer of resistance $50\unit{\Omega}$ is connected to a battery of $3\unit{V}$ along with a resistance of $2950\unit{\Omega}$ in series shows full-scale deflection of 30 divisions. The additional series resistance required to reduce the deflection to 20 divisions is}
{\opt{$1500\unit{\Omega}$}{$4440\unit{\Omega}$}{$7400\unit{\Omega}$}{$2950\unit{\Omega}$}}

\QQ{The magnetic field at the center of a current carrying loop of radius $0.1\unit{m}$ is $5\sqrt5$ times that at a point along its axis. The distance of this point from the centre of the loop is}
{\opt{$0.2\unit{m}$}{$0.1\unit{m}$}{$0.05\unit{m}$}{$0.25\unit{m}$}}

\QQ{A straight wire of length $50\unit{cm}$ carrying a current of $2.5\unit{A}$ is suspended in mid-air by a uniform magnetic field of $0.5\unit{T}$ (as shown in figure). The mass of the wire is ($g = 10\unit{m\,s^{-2}}$)
\fig{\begin{tikzpicture}[scale=0.8]
\draw[->] (0,0) -- (3,0) node[midway,below,lbl]{$I$}; \node[lbl] at (2.2,-0.6) {$\otimes$ B};
\end{tikzpicture}}}
{\opt{$62.5\unit{g}$}{$250\unit{g}$}{$125\unit{g}$}{$100\unit{g}$}}

\QQ{Which of the following properties is `False' for a bar magnet?}
{\opt{Its poles cannot be separated.}{It points in North-South direction when suspended.}{Its like poles repel and unlike poles attract.}{It doesn't produce magnetic field.}}

\QQ{A magnetic dipole of magnetic moment $6\times10^{-2}\unit{A\,m^2}$ and moment of inertia $12\times10^{-6}\unit{kg\,m^2}$ performs oscillations in a magnetic field of $2\times10^{-2}\unit{T}$. The time taken by the dipole to complete 20 oscillations is ($\pi\simeq3$)}
{\opt{$36\unit{s}$}{$6\unit{s}$}{$12\unit{s}$}{$18\unit{s}$}}

\QQ{The susceptibility of a ferromagnetic substance is}
{\opt{$\gg1$}{$>1$}{$<1$}{Zero}}

\QQ{A bar magnet is allowed to fall vertically through a copper coil placed in a horizontal plane. The magnet falls with a net acceleration
\fig{\begin{tikzpicture}[scale=0.7]
\draw (-0.25,1.0) rectangle (0.25,2.4); \node[lbl] at (0,2.1) {S}; \node[lbl] at (0,1.3) {N};
\draw[densely dashed,->] (0,1.0) -- (0,0.05); \draw (0,0) ellipse (1.2 and 0.45);
\end{tikzpicture}}}
{\opt{$= g$}{$> g$}{$< g$}{Zero}}

\QQ{The working of magnetic braking of trains is based on}
{\opt{Alternating current}{Eddy current}{Steady current}{Pulsating current}}

\QQ{A jet plane of wing span $20\unit{m}$ is travelling towards west at a speed of $400\unit{m\,s^{-1}}$. If the earth's total magnetic field is $4\times10^{-4}\unit{T}$ and the dip angle is $30^\circ$, at that place, the voltage difference developed across the ends of the wing is}
{\opt{$1.6\unit{V}$}{$3.2\unit{V}$}{$0.8\unit{V}$}{$6.4\unit{V}$}}

\QQ{In the A.C.\ circuit shown, keeping `K' pressed, if an iron rod is inserted into the coil, the bulb in the circuit,
\fig{\begin{circuitikz}[scale=0.6,transform shape,font=\Large]
\draw (0,0) to[sV,l_=A.C] (0,2.2) to[lamp,l=Bulb] (2.4,2.2) to[L,l=Coil] (5,2.2) -- (6,2.2) -- (6,0) to[push button,l_=K] (3.6,0) -- (0,0);
\fill[pattern=north east lines] (5.4,2.6) rectangle (7.2,2.85); \node at (6.3,3.2) {Rod};
\end{circuitikz}}}
{\opt{glows more brightly}{glows less brightly}{glows with same brightness (as before the rod is inserted)}{gets damaged}}

\QQ{The output of a step down transformer is measured to be $48\unit{V}$ when connected to a $12\unit{W}$ bulb. The value of peak current is}
{\opt{$\dfrac1{\sqrt2}\unit{A}$}{$\sqrt2\unit{A}$}{$\dfrac1{2\sqrt2}\unit{A}$}{$\dfrac14\unit{A}$}}

\QQ{A coil of inductive reactance $1/\sqrt3\unit{\Omega}$ and resistance $1\unit{\Omega}$ is connected to a $200\unit{V}$, $50\unit{Hz}$ A.C.\ supply. The time lag between maximum voltage and current is}
{\opt{$\dfrac1{300}\unit{s}$}{$\dfrac1{600}\unit{s}$}{$\dfrac1{500}\unit{s}$}{$\dfrac1{200}\unit{s}$}}

\QQ{If $\vec E$ and $\vec B$ represent electric and magnetic field vectors of an electromagnetic wave, the direction of propagation of the wave is along}
{\opt{$\vec E$}{$\vec B$}{$\vec E\times\vec B$}{$\vec B\times\vec E$}}

\QQ{According to Cartesian sign convention, in ray optics}
{\opt{all distances are taken positive}{all distances are taken negative}{all distances in the direction of incident ray are taken positive}{all distances in the direction of incident ray are taken negative}}

\QQ{A linear object of height $10\unit{cm}$ is kept in front of a concave mirror of radius of curvature $15\unit{cm}$, at a distance of $10\unit{cm}$. The image formed is}
{\opt{magnified and erect}{magnified and inverted}{diminished and erect}{diminished and inverted}}

\QQ{During scattering of light, the amount of scattering is inversely proportional to \rule{1cm}{0.4pt} of wavelength of light.}
{\opt{cube}{square}{fourth power}{half}}

\QQ{In Young's double-slit experiment if yellow light is replaced by blue light, the interference fringes become}
{\opt{wider}{narrower}{brighter}{darker}}

\QQ{According to Huygens' principle, during refraction of light from air to a denser medium}
{\opt{Wavelength and speed decrease}{Wavelength and speed increase}{Wavelength increases but speed decreases}{Wavelength decreases but speed increases}}

\QQ{In a system of two crossed polarisers, it is found that the intensity of light from the second polariser is half from that of first polariser. The angle between their pass axes is}
{\opt{$45^\circ$}{$60^\circ$}{$30^\circ$}{$0^\circ$}}

\QQ{From the following graph of photo current against collector plate potential, for two different intensities of light $I_1$ and $I_2$, one can conclude
\fig{\begin{tikzpicture}[scale=0.8]
\draw[->] (-1.6,0) -- (2.2,0) node[below left,lbl]{Collector plate potential}; \draw[->] (0,0) -- (0,1.9) node[above,lbl]{Photo current};
\draw (-1.4,0) .. controls (-0.6,0.5) and (-0.2,1.4) .. (0.6,1.45) -- (1.8,1.45) node[right,lbl]{$I_2$};
\draw (-1.4,0) .. controls (-0.6,0.4) and (-0.2,1.0) .. (0.6,1.05) -- (1.8,1.05) node[right,lbl]{$I_1$};
\node[lbl,below] at (-1.4,0) {$-V_0$};
\end{tikzpicture}}}
{\opt{$I_1 = I_2$}{$I_1>I_2$}{$I_1<I_2$}{Comparison is not possible.}}

\QQ{A particle is dropped from a height `$H$'. The de Broglie wavelength of the particle depends on height as}
{\opt{$H$}{$H^0$}{$H^{1/2}$}{$H^{-1/2}$}}

\QQ{The scientist who is credited with the discovery of `nucleus' in an atom is}
{\opt{J.J.\ Thomson}{Rutherford}{Niels Bohr}{Balmer}}

\QQ{The energy (in eV) required to excite an electron from $n = 2$ to $n = 4$ state in hydrogen atom is}
{\opt{$+2.55$}{$-3.4$}{$-0.85$}{$+4.25$}}

\QQ{In a nuclear reactor the function of the Moderator is to decrease}
{\opt{number of neutrons}{speed of neutrons}{escape of neutrons}{temperature of the reactor}}

\QQ{The particles emitted in the decay of ${}^{238}_{92}\mathrm U$ to ${}^{234}_{92}\mathrm U$}
{\opt{1 $\alpha$ and 2 $\beta$}{1 $\alpha$ only}{1 $\alpha$ and 1 $\beta$}{2 $\alpha$ and 2 $\beta$}}

\QQ{The mass defect of ${}^4_2\mathrm{He}$ is $0.03\unit{u}$. The binding energy per nucleon of helium (in MeV) is}
{\opt{27.93}{6.9825}{2.793}{69.825}}

\QQ{The energy gap in case of which of the following is less than $3\unit{eV}$?}
{\opt{Copper}{Iron}{Aluminium}{Germanium}}

\QQ{Which of the following semi-conducting devices is used as voltage regulator?}
{\opt{Photo diode}{LASER diode}{Zener diode}{Solar cell}}

\QQ{In the three parts of a transistor, `Emitter' is of}
{\opt{moderate size and heavily doped}{large size and lightly doped}{thin size and heavily doped}{large size and moderately doped}}

\QQ{In the figure shown, if the diode forward voltage drop is $0.2\unit{V}$, the voltage difference between A and B is
\fig{\begin{circuitikz}[scale=0.6,transform shape,font=\Large]
\draw (0,5) node[circ]{} node[left]{A} to[short,i=$0.2\,\mathrm{mA}$] (0,4.2) to[R,l=$5\,\mathrm{k\Omega}$] (0,2.8) -- (0,2.6) to[empty diode] (0,1.6) to[R,l=$5\,\mathrm{k\Omega}$] (0,0) node[circ]{} node[left]{B};
\draw (0,2.7) to[C] (2,2.7);
\end{circuitikz}}}
{\opt{$1.3\unit{V}$}{$2.2\unit{V}$}{0}{$0.5\unit{V}$}}

\QQ{Which of the following logic gates is considered as `universal'?}
{\optfig{\begin{circuitikz}[scale=0.75,transform shape]\draw (0,0) node[xor port](g){}; \draw (g.in 1) node[left]{A}; \draw (g.in 2) node[left]{B}; \draw (g.out) node[right]{Y};\end{circuitikz}}
{\begin{circuitikz}[scale=0.75,transform shape]\draw (0,0) node[left]{A} to[short] (0.3,0) node[buffer port,anchor=in](g){}; \draw (g.out) node[right]{Y};\end{circuitikz}}
{\begin{circuitikz}[scale=0.75,transform shape]\draw (0,0) node[and port](g){}; \draw (g.in 1) node[left]{A}; \draw (g.in 2) node[left]{B}; \draw (g.out) node[right]{Y};\end{circuitikz}}
{\begin{circuitikz}[scale=0.75,transform shape]\draw (0,0) node[nand port](g){}; \draw (g.in 1) node[left]{A}; \draw (g.in 2) node[left]{B}; \draw (g.out) node[right]{Y};\end{circuitikz}}}

\QQ{A basic communication system consists of
\begin{stmts}[(\alph*)]
\item Transmitter
\item Information source
\item User of information
\item Channel
\item Receiver
\end{stmts}
The correct sequence of the arrangement is}
{\opt{a, b, c, d and e}{b, a, d, e and c}{b, d, a, c and e}{b, e, a, d and c}}

\end{multicols}
